\chapter{Why Strings, Why Two Scales, Why a Proof Assistant}\label{ch:intro}

This chapter answers three questions before any machinery is built. Why would one replace
point particles by strings? Why does a string, unlike a particle, treat a circle of radius $R$
and a circle of radius $\ap/R$ as the same space, so that every length comes with a partner?
And why does this book accompany the physics with statements checked by a proof assistant?
A student should care about the first two because they contain, in a form that needs only
quantum mechanics and special relativity, the idea that the rest of the book develops with
increasing generality: from one circle to tori, to the doubled geometry of double field
theory, to K3 surfaces and their lattices. A student should care about the third because
the subject mixes theorems, established physics and open conjectures, and it is easy to lose
track of which is which. The chapter ends with one small result carried through every
stage: derived by hand, checked with numbers, and stated as a theorem of Lean~4.

%=====================================================================================
\section{Why strings}\label{sec:intro:whystrings}

\subsection{Gravity resists the usual quantization}

Units with $\hbar=c=1$ are used throughout, so that energy, mass and inverse length are the
same dimension. Newton's constant $G_N$ then defines a mass, the \emph{Planck mass}%
\index{Planck mass}. In the normalization of the Einstein--Hilbert action
$S_{EH}=\frac{1}{16\pi G_N}\int \dd^4x\,\sqrt{-g}\,\mathcal{R}$ one writes
\begin{equation}\label{eq:intro:planck}
  8\pi G_N=\frac{1}{M_{\rm pl}^2},\qquad M_{\rm pl}\approx 2\times 10^{18}\ \text{GeV},
\end{equation}
the \emph{reduced} Planck mass; without the $8\pi$ one finds $1/\sqrt{G_N}\approx 10^{19}$ GeV
\source{Tong §0.1, ll.~465--493}. Here $g$ is the determinant of the spacetime metric
$g_{\mu\nu}$ and $\mathcal{R}$ its Ricci scalar.

To see what \cref{eq:intro:planck} means for a quantum theory, expand the metric around flat
space, $g_{\mu\nu}=\eta_{\mu\nu}+h_{\mu\nu}/M_{\rm pl}$. The factor $1/M_{\rm pl}$ is chosen so
that the kinetic term of the fluctuation $h$ carries no power of $M_{\rm pl}$. The action then
has the schematic form
\begin{equation}\label{eq:intro:expansion}
  S_{EH}=\int \dd^4x\,\Big[(\partial h)^2+\frac{1}{M_{\rm pl}}\,h(\partial h)^2
  +\frac{1}{M_{\rm pl}^2}\,h^2(\partial h)^2+\dots\Big]
\end{equation}
\source{Tong §0.1, ll.~494--523}. Every interaction is suppressed by a power of $M_{\rm pl}$.
A process at energy $E$ is therefore controlled by the dimensionless ratio
$E^2/M_{\rm pl}^2$: gravity is weak at low energy, and perturbation theory fails when $E$
approaches $M_{\rm pl}$. A coupling with the dimension of an inverse mass is the signature of
a \emph{non-renormalizable}\index{non-renormalizability} theory: the divergences of its loop
diagrams cannot be absorbed in a finite number of counterterms. For pure gravity the first
such divergence appears at two loops; with matter, at one loop
\source{Tong §0.1, ll.~531--560}.

This does not make general relativity useless. It makes it an \emph{effective field theory},
accurate to order $(E/M_{\rm pl})^2$, in the same way that Fermi's four-fermion theory of the
weak interaction is accurate below about $100$ GeV. Fermi's theory fails at its cutoff because
it lacks the degrees of freedom that tame the high-energy behaviour, the $W$ and $Z$ bosons.
The question of quantum gravity, put in this language, is: what are the missing degrees of
freedom of gravity? \source{Tong §0.1, ll.~572--587}

\subsection{The string hypothesis and its one parameter}

String theory answers with a change of the elementary object: matter consists of tiny loops
of string, not of points \source{Tong §0, ll.~373--378}. A relativistic string is
characterized by a single dimensionful number, its tension\index{string tension} $T$, the
energy per unit length: a static string of length $L$ has potential energy $TL$, growing
linearly with $L$ (unlike a Hooke spring, where it grows quadratically)
\source{Tong §1.2, ll.~1090--1096}. For historical reasons the tension is written
\begin{equation}\label{eq:intro:tension}
  T=\frac{1}{2\pi\ap},\qquad \ap=\ell_s^2 ,
\end{equation}
where $\ap$ is the \emph{Regge slope}\index{Regge slope} and $\ell_s$ the \emph{string
length}\index{string length} \source{Tong §1.2, eqs.~(1.17)--(1.18), ll.~1097--1120}.
A dimensional check: $T$ is energy per length, that is, mass squared in natural units, so
$\ap$ is a length squared and $\ell_s=\sqrt{\ap}$ is a length. These are the conventions of the
whole book (mostly-plus signature, $\ap$ of dimension length$^2$).

For the fundamental strings of this book, which have no thickness and no further structure,
the tension is taken far above any laboratory scale, ``typically an order of magnitude or so
below the Planck scale'' \source{Tong §1.2, ll.~1146--1150}. To convert to a length one uses
$\hbar c\approx 1.97\times10^{-14}$ GeV\,cm (a standard value; not checked against a source in
this book's library): an energy of $10^{19}$ GeV corresponds to
$1.97\times 10^{-14}/10^{19}\approx 2\times10^{-33}$ cm, in agreement with the figure
$10^{-33}$ cm quoted for the Planck length in \source{GPR §1, ll.~209--213}. By the same
conversion the Fermi scale of $100$ GeV is $2\times10^{-16}$ cm. Seventeen orders of magnitude
separate the two, and this is the practical reason why no experiment tests the ideas of this
book directly. Tong says so without hedging: string theory is ``speculative science'', with no
experimental evidence that it describes our world \source{Tong §0, ll.~380--387}.

\begin{physicsbox}[Physical picture: what the new length does]
A point particle can probe arbitrarily short distances if it is given enough energy. A string
cannot: it is itself an extended, fluctuating object of size about $\ell_s$, so that spacetime
looks ``fuzzy'' to it at the scale $\sqrt{\ap}$. One consequence is that string scattering
amplitudes fall off exponentially at high energy, where those of general relativity grow; this
is the first sign that the ultraviolet problem of \cref{eq:intro:expansion} may be cured
\source{Tong §6, ll.~8055--8080}. The string length is a \emph{minimum length}%
\index{minimum length}, but ``not in any crude simple manner'': it shows up differently
depending on the question asked. \Cref{sec:intro:twoscales} describes the manifestation that
organizes this book.
\end{physicsbox}

Why study such a theory? The source gives four reasons that a student should weigh
independently \source{Tong §0, ll.~388--463}: it is \emph{a} consistent framework for quantum
gravity, in which sharp questions about short-distance spacetime can be posed; it may be
\emph{the} theory of quantum gravity, since general relativity, gauge theories, scalars and
chiral fermions all emerge at low energy; it gives new tools for ordinary gauge theories; and
it produces new mathematics, mirror symmetry being the best-known case. This book lives
mostly in the fourth reason. Its subject is a part of string theory where the mathematics is
rigid enough --- lattices, integer matrices, modular forms, finite groups --- for a proof
assistant to check it.

\begin{historybox}
String theory ``was born from attempts to understand the strong force''
\source{Tong §0, ll.~437--438}; the name ``Regge slope'' for $\ap$ is a remnant of that
period \source{Tong §3.1.3, ll.~3391--3394}. The word \emph{duality} has a longer record in
physics: wave--particle duality, the high/low temperature duality of the Ising model, and
electric--magnetic duality all precede the target-space duality of strings; according to
the standard review each shares features with it and ``none is identical to it''
\source{GPR §1, ll.~237--255}.
\end{historybox}

%=====================================================================================
\section{Why two scales}\label{sec:intro:twoscales}

The cleanest way to see what is new about a string is to put it on the simplest space that is
not $\R^n$: one spatial direction curled into a circle,
\begin{equation}\label{eq:intro:circle}
  X\sim X+2\pi R ,
\end{equation}
with $R$ the radius. \Cref{ch:circle,ch:tduality} treat this carefully; here only the two
facts that need no string quantization are derived.

\subsection{A particle on a circle sees one scale}

A quantum particle with momentum $p$ along the circle has the wavefunction
$\psi(X)=\ee^{\ii pX}$. The points $X$ and $X+2\pi R$ are the same point, so $\psi$ must take
the same value there: $\ee^{2\pi\ii pR}=1$, hence
\begin{equation}\label{eq:intro:momentum}
  p=\frac{n}{R},\qquad n\in\Z
\end{equation}
\source{Tong §8.2, ll.~11352--11358}. Suppose the particle is massless in the full spacetime,
$-p_\mu p^\mu-p^2=0$, where $\mu$ runs over the non-compact directions. An observer who
cannot resolve the circle measures the mass $M^2=-p_\mu p^\mu$, and finds a tower
\begin{equation}\label{eq:intro:kk}
  M^2_{\text{particle}}=\frac{n^2}{R^2},\qquad n\in\Z ,
\end{equation}
the \emph{Kaluza--Klein tower}\index{Kaluza--Klein tower}. The spacing $1/R$ of the tower
measures the radius: a small circle gives heavy states, a large circle gives light ones, and
two different radii are never confused. As $R\to0$ every state with $n\neq0$ becomes
infinitely heavy and the circle direction disappears from low-energy physics.

\subsection{A string on a circle sees two}

A closed string can do what a particle does: move around the circle with momentum
\eqref{eq:intro:momentum}. It can also do something a particle cannot: wrap around the circle.
Parametrize the string by $\sigma\in[0,2\pi]$. Closing up after one turn of $\sigma$ only
requires returning to the same \emph{point of the circle}, so by \eqref{eq:intro:circle}
\begin{equation}\label{eq:intro:winding}
  X(\sigma+2\pi)=X(\sigma)+2\pi wR,\qquad w\in\Z ,
\end{equation}
where $w$ is the \emph{winding number}\index{winding number}
\source{Tong §8.2, ll.~11359--11371; Tong writes $m$ for our $w$}. A string wound $w$ times has
length at least $2\pi|w|R$, and a string costs energy $T$ per unit length. With
\eqref{eq:intro:tension} the winding contributes
\begin{equation}\label{eq:intro:windingmass}
  M_{\text{winding}}=2\pi |w| R\,T=\frac{|w|R}{\ap}
\end{equation}
to the mass \source{Tong §8.2, ll.~11480--11487}. \Cref{fig:intro:cylinder} shows the two
kinds of excitation.

\begin{figure}[t]
\centering
\begin{tikzpicture}[scale=0.95]
  % left cylinder: momentum
  \begin{scope}
    \draw (0,0) ellipse (0.45 and 1.1);
    \draw (0,1.1) -- (4.2,1.1);
    \draw (0,-1.1) -- (4.2,-1.1);
    \draw (4.2,1.1) arc (90:-90:0.45 and 1.1);
    \draw[dashed] (4.2,1.1) arc (90:270:0.45 and 1.1);
    % small closed string moving around
    \draw[thick,blue!60!black] (2.1,0.15) ellipse (0.33 and 0.2);
    \draw[-{Stealth},blue!60!black] (2.1,0.45) -- (2.1,0.95);
    \node[blue!60!black] at (2.95,0.72) {\small $p=n/R$};
    \node at (2.1,-1.6) {\small momentum: $M\sim |n|/R$};
  \end{scope}
  % right cylinder: winding
  \begin{scope}[xshift=6.6cm]
    \draw (0,0) ellipse (0.45 and 1.1);
    \draw (0,1.1) -- (4.2,1.1);
    \draw (0,-1.1) -- (4.2,-1.1);
    \draw (4.2,1.1) arc (90:-90:0.45 and 1.1);
    \draw[dashed] (4.2,1.1) arc (90:270:0.45 and 1.1);
    % wound string
    \draw[thick,red!70!black] (2.1,1.1) arc (90:-90:0.45 and 1.1);
    \draw[thick,red!70!black,dashed] (2.1,1.1) arc (90:270:0.45 and 1.1);
    \node[red!70!black] at (3.25,0) {\small $w=1$};
    \node at (2.1,-1.6) {\small winding: $M\sim |w|R/\ap$};
    \draw[{Stealth}-{Stealth}] (-0.85,-1.1) -- node[left]{\small $2\pi R$} (-0.85,1.1);
  \end{scope}
\end{tikzpicture}
\caption{One compact direction, drawn as the circumference of a cylinder; the axis of the
cylinder stands for the non-compact directions. Left: a small closed string moving around
the circle carries quantized momentum $n/R$, as a particle would. Right: a string wrapped
around the circle carries winding number $w$; its minimal length $2\pi|w|R$ costs the energy
\eqref{eq:intro:windingmass}. A particle has no analogue of the second state.}
\label{fig:intro:cylinder}
\end{figure}

Quantizing the string (\cref{ch:quantum,ch:circle}) adds the oscillator levels $N,\tilde N\in
\{0,1,2,\dots\}$ of the right- and left-moving vibrations and a zero-point term. The result,
for the bosonic string and in the conventions of this book, is
\begin{equation}\label{eq:intro:mass}
  M^2=\frac{n^2}{R^2}+\frac{w^2R^2}{\ap^2}+\frac{2}{\ap}\big(N+\tilde N-2\big),
  \qquad N-\tilde N=nw ,
\end{equation}
\source{Tong §8.2, eqs.~(8.5)--(8.6), ll.~11425--11475}. The first term is the
particle's \eqref{eq:intro:kk}; the second is the square of \eqref{eq:intro:windingmass}; the
third is the vibrational energy, in units set by the tension. The condition $N-\tilde N=nw$ is
\emph{level matching}\index{level matching}: it expresses that no point of a closed string is
preferred as the origin of $\sigma$. \Cref{eq:intro:mass} can be packaged with the
combinations $p_{L,R}=n/R\pm wR/\ap$: a line of algebra gives
$\tfrac12(p_L^2+p_R^2)=n^2/R^2+w^2R^2/\ap^2$ and $p_L^2-p_R^2=4nw/\ap$, so that both the mass
and the level-matching condition are functions of $(p_L,p_R)$ alone. This is the seed of the
Narain lattice of \cref{ch:narain}.

Now compare a string on a circle of radius $R$ with a string on a circle of radius
\begin{equation}\label{eq:intro:tdual}
  \tilde R=\frac{\ap}{R}.
\end{equation}
Substituting in \eqref{eq:intro:mass}, $n^2/\tilde R^2=n^2R^2/\ap^2$ and
$w^2\tilde R^2/\ap^2=w^2/R^2$: the first two terms are exchanged. If the labels are exchanged
as well, $n\leftrightarrow w$, the formula returns to itself, and level matching, which
contains only the product $nw$, is untouched. Hence
\begin{equation}\label{eq:intro:tduality}
  M^2(R;\,n,w,N,\tilde N)=M^2\!\big(\ap/R;\,w,n,N,\tilde N\big):
\end{equation}
the two spectra coincide state by state \source{Tong §8.3, eqs.~(8.7)--(8.8),
ll.~11626--11640; BW §III.B, eq.~(24), ll.~720--737}. This is
\emph{T-duality}\index{T-duality}. The statement extends from the spectrum to the full
conformal field theory and hence to interactions \source{Tong §8.3, ll.~11655--11658} \TL;
\cref{ch:tduality,ch:buscher} explain why. The standard review puts it this way: a string
``cannot tell'' whether it propagates on a circle of radius $R$ or $1/R$ in string units
\source{GPR §1, ll.~229--234}.

\begin{physicsbox}[Physical picture: two rulers]
A particle owns one ruler: the spacing $1/R$ of its momentum tower. A string owns two: the
momentum spacing $1/R$ and the winding spacing $R/\ap$. When $R\gg\sqrt{\ap}$ the momentum
states are light and nearly continuous --- which is what ``a large dimension'' means
operationally --- and the winding states are too heavy to matter. When $R\ll\sqrt{\ap}$ the
roles are reversed: the \emph{winding} states form a near-continuum, and the spectrum looks
``as if another dimension of space is opening up'' \source{Tong §8.3, ll.~11641--11654}. An
observer made of strings who measures the size of the circle by the lightest tower reads off
$\max(R,\ap/R)$, never less than $\sqrt{\ap}$. This is the sense in which the theory has two
scales, $R$ and $\ap/R$, of equal standing.
\end{physicsbox}

The fixed point of \eqref{eq:intro:tdual}, $R=\sqrt{\ap}=\ell_s$, is the \emph{self-dual
radius}\index{self-dual radius}. Shrinking a circle below it is the same as growing the dual
circle, so $\sqrt{\ap}$ acts as a minimal radius \source{Tong §8.3, ll.~11659--11663}.
\Cref{fig:intro:towers} plots the two rulers and their sum.

\begin{figure}[t]
\centering
\begin{tikzpicture}[xscale=2.0,yscale=0.95]
  \draw[-{Stealth}] (0,0) -- (3.55,0) node[right] {\small $R/\sqrt{\ap}$};
  \draw[-{Stealth}] (0,0) -- (0,4.4) node[above] {\small mass or length, in units of $\sqrt{\ap}$};
  \foreach \x in {1,2,3} \draw (\x,0.06) -- (\x,-0.06) node[below] {\small $\x$};
  \foreach \y in {1,2,3,4} \draw (0.03,\y) -- (-0.03,\y) node[left] {\small $\y$};
  \draw[thick,blue!60!black,domain=0.245:3.4,samples=80] plot (\x,{1/\x});
  \draw[thick,red!70!black,domain=0:3.4] plot (\x,{\x});
  \draw[thick,domain=0.27:3.4,samples=80] plot (\x,{\x+1/\x});
  \draw[dashed] (0,2) -- (3.4,2);
  \draw[dotted] (1,0) -- (1,2);
  \fill (1,2) circle (0.9pt);
  \fill (1,1) circle (0.9pt);
  \node[blue!60!black] at (0.62,3.3) {\small $1/R$};
  \node[red!70!black] at (3.1,2.6) {\small $R/\ap$};
  \node at (2.2,3.55) {\small $R+\ap/R$};
  \node at (2.75,1.78) {\small bound $2\sqrt{\ap}$};
\end{tikzpicture}
\caption{The two rulers of a string on a circle, in units $\sqrt{\ap}=1$: the mass $1/R$ of
the first momentum state (blue) and the mass $R/\ap$ of the first winding state (red) cross at
the self-dual radius $R=\sqrt{\ap}$. The black curve is the effective scale $R+\ap/R$ of
\cref{sec:intro:amgm}; it is symmetric under $R\mapsto\ap/R$ and never drops below
$2\sqrt{\ap}$ (dashed).}
\label{fig:intro:towers}
\end{figure}

\begin{example}[The same spectrum at $R=2$ and $R=\tfrac12$]\label{ex:intro:spectrum}
Set $\ap=1$, so that lengths are measured in string lengths, and compare $R=2$ with its dual
$\tilde R=\tfrac12$. \Cref{eq:intro:mass} reads $M^2=n^2/R^2+w^2R^2+2(N+\tilde N-2)$.
Each row of the table is a state allowed by level matching ($N-\tilde N=nw$ is checked in the
second column).
\begin{center}
\begin{tabular}{@{}ccccc@{}}
\toprule
$(n,w)$ & $(N,\tilde N)$ & $M^2$ at $R=2$ & $M^2$ at $R=\tfrac12$ & particle, $n^2/R^2$ at
$R=2$ / $R=\tfrac12$\\
\midrule
$(0,0)$ & $(1,1)$ & $0$ & $0$ & $0$ / $0$\\
$(1,0)$ & $(1,1)$ & $\tfrac14$ & $4$ & $\tfrac14$ / $4$\\
$(0,1)$ & $(1,1)$ & $4$ & $\tfrac14$ & --- \\
$(2,0)$ & $(1,1)$ & $1$ & $16$ & $1$ / $16$\\
$(0,2)$ & $(1,1)$ & $16$ & $1$ & --- \\
$(1,1)$ & $(1,0)$ & $\tfrac94$ & $\tfrac94$ & --- \\
\bottomrule
\end{tabular}
\end{center}
For instance, for $(n,w)=(1,1)$, $(N,\tilde N)=(1,0)$ at $R=2$:
$M^2=\tfrac14+4+2(1+0-2)=\tfrac94$. Read as \emph{sets}, the third and fourth columns agree:
$\{0,\tfrac14,4,1,16,\tfrac94\}$ in both cases. Row by row they agree after the exchange
$n\leftrightarrow w$: the state $(1,0)$ at $R=2$ has the mass of the state $(0,1)$ at
$R=\tfrac12$. The last column shows what a particle sees: the towers $\{0,\tfrac14,1,\dots\}$ and
$\{0,4,16,\dots\}$ are different, so a particle distinguishes the two circles and a string
does not. (The level $N=\tilde N=0$, $n=w=0$ gives $M^2=-4/\ap$, the tachyon of the bosonic
string; it is discussed in \cref{ch:quantum} and plays no role here.)
\end{example}

\begin{pitfallbox}
T-duality does \emph{not} say that the function $M^2(R)$ of a given state is symmetric under
$R\mapsto\ap/R$; the table shows that it is not. It says that the two \emph{theories} have the
same set of states, under a relabelling that exchanges momentum and winding. A second trap:
the conclusion depends on having both towers. Any approximation that keeps only the momentum
states --- ordinary field theory on the circle, for example --- loses the duality. Double field
theory (\cref{ch:doubled}) is the attempt to write a field theory that keeps both.
\end{pitfallbox}

%=====================================================================================
\section{A toy theorem, three ways: \texorpdfstring{$R+\ap/R\ge 2\sqrt{\ap}$}{R + a'/R >= 2 sqrt a'}}
\label{sec:intro:amgm}

\subsection{By hand}

A quantity that treats $R$ and $\ap/R$ on an equal footing must be invariant under the
exchange \eqref{eq:intro:tdual}. The simplest one with the dimension of a length is the sum
\begin{equation}\label{eq:intro:reff}
  R_{\rm eff}(R)=R+\frac{\ap}{R},\qquad R>0 ,
\end{equation}
which this book calls the \emph{effective scale}\index{effective scale}\index{dual-scale
principle}. Invariance is immediate: $R_{\rm eff}(\ap/R)=\ap/R+\ap/(\ap/R)=\ap/R+R$.

\begin{proposition}\label{thm:intro:amgm}
For all real $R>0$ and $\ap>0$, $R+\ap/R\ge2\sqrt{\ap}$, with equality if and only if
$R=\sqrt{\ap}$.
\end{proposition}

\begin{proof}[First proof (completing the square)]
Since $R>0$ we may write $R=(\sqrt R)^2$ and $\ap/R=(\sqrt{\ap}/\sqrt R)^2$, and the product of
the two square roots is $\sqrt{\ap}$. Therefore
\begin{equation}\label{eq:intro:square}
  R+\frac{\ap}{R}-2\sqrt{\ap}=\Big(\sqrt R-\frac{\sqrt{\ap}}{\sqrt R}\Big)^{2}
  =\frac{(R-\sqrt{\ap})^2}{R}\;\ge 0 .
\end{equation}
The middle expression is a square, the right-hand one is a square divided by a positive
number; either way it is non-negative and vanishes exactly when $R=\sqrt{\ap}$.
\end{proof}

\begin{proof}[Second proof (calculus)]
$R_{\rm eff}'(R)=1-\ap/R^2$ vanishes on $R>0$ only at $R=\sqrt{\ap}$, and
$R_{\rm eff}''(R)=2\ap/R^3>0$, so $R_{\rm eff}$ is convex with a single minimum,
$R_{\rm eff}(\sqrt{\ap})=2\sqrt{\ap}$.
\end{proof}

\begin{proof}[Third proof (arithmetic and geometric means)]
For $a,b\ge0$, $(a+b)/2\ge\sqrt{ab}$, the AM--GM inequality\index{AM--GM inequality}, which is
itself $(\sqrt a-\sqrt b)^2\ge0$. With $a=R$ and $b=\ap/R$ the geometric mean is
$\sqrt{R\cdot\ap/R}=\sqrt{\ap}$, independent of $R$: \emph{T-dual radii have a fixed geometric
mean, so their arithmetic mean is bounded below by it.}
\end{proof}

The identities in \eqref{eq:intro:square} were also confirmed with a computer-algebra
system (a symbolic check, which is not a proof-assistant check). It is useful to remove the
units. Put $x=R/\sqrt{\ap}$, the radius in string lengths. Then
\begin{equation}\label{eq:intro:dimless}
  R+\frac{\ap}{R}=\sqrt{\ap}\,\Big(x+\frac1x\Big),\qquad\text{so}\qquad
  R+\frac{\ap}{R}\ge2\sqrt{\ap}\iff x+\frac1x\ge2 .
\end{equation}
The dimensionless statement $x+1/x\ge2$ for $x>0$ is the one that appears in Lean below.

\begin{example}[Numbers]\label{ex:intro:numbers}
In units $\sqrt{\ap}=1$: at $R=2$, $R_{\rm eff}=2+\tfrac12=2.5$; at the dual radius
$R=\tfrac12$, $R_{\rm eff}=\tfrac12+2=2.5$ again; at $R=1$, $R_{\rm eff}=2$, the minimum. The
excess over the bound, from \eqref{eq:intro:square}, is $(R-1)^2/R$: for $R=2$ this is
$\tfrac12$, for $R=\tfrac12$ it is $(\tfrac14)/(\tfrac12)=\tfrac12$ as well, as it must be. For a
circle ten times the string length, $R_{\rm eff}=10.1$: the correction to the naive size is
one per cent, and for a macroscopic circle it is unobservably small, $\ap/R^2$ in relative
terms. The effective scale differs from $R$ only near or below the string length.
\end{example}

\subsection{In Lean}

The same statement, in the dimensionless form \eqref{eq:intro:dimless}, is a theorem of this
book's Lean library. It is the first \texttt{leanbox} of the book, so it is annotated in more
detail than later ones.

\begin{leanbox}[In Lean: the circle bound \TA]
\begin{leancode}
-- DualScaleStream2/DualScale/TraceBound.lean
theorem circle_effective_scale_ge_two (R : ℝ) (hR : 0 < R) :
    2 ≤ R + R⁻¹ := by
  have h : R + R⁻¹ - 2 = (R - 1) ^ 2 / R := by field_simp; ring
  have h_nonneg : 0 ≤ (R - 1) ^ 2 / R := by
    apply div_nonneg
    · exact sq_nonneg (R - 1)
    · linarith
  linarith
\end{leancode}
\textbf{Reading.} For every real number $R$ and every proof \lean{hR} that $0<R$, one has
$2\le R+R^{-1}$. Nothing in the statement mentions strings, $\ap$ or a circle: $R$ is a real
number, and the identification $R=(\text{radius})/\sqrt{\ap}$ is made by the reader.

\textbf{Proof idea.} Exactly the first proof above with $\ap=1$: the tactic \lean{field\_simp}
clears the denominator and \lean{ring} verifies the polynomial identity
$R+R^{-1}-2=(R-1)^2/R$; a square divided by a positive number is non-negative
(\lean{div\_nonneg}, \lean{sq\_nonneg}); \lean{linarith} combines the two facts.
\end{leanbox}

\paragraph{How to read a \texttt{leanbox}.} Every box of this kind in the book has the same
four parts.
\begin{enumerate}[leftmargin=*,itemsep=2pt]
\item \emph{The statement, verbatim.} Text before the colon declares the data and the
hypotheses; a hypothesis is a named proof, here \lean{hR : 0 < R}. Text after the colon is
the conclusion. The part after \texttt{:=} is the proof; from now on it is usually abbreviated
to \texttt{:= by ...}, because for the reader the \emph{statement} is what matters --- the
proof has been checked by the Lean kernel and need not be trusted or even read.
\item \emph{A plain-language reading} that says what the statement literally asserts and no
more.
\item \emph{The proof idea}, two or three sentences, so that the result can be connected to
the derivation in the text.
\item \emph{A tier} (\cref{sec:intro:tiers}): \TA{} for the Lean statement, and separately the
tier of the physical statement that it supports.
\end{enumerate}
The gap between the Lean statement and \cref{thm:intro:amgm} is small but real, and is typical.
The Lean theorem is the case $\ap=1$; the general case follows from the rescaling
\eqref{eq:intro:dimless}, which was done on this page and not in Lean. The equality case
($R=1$ is the only minimizer) is not part of the Lean statement either; \cref{exr:intro:eq}
asks for it.

\begin{pitfallbox}[Common pitfall: the hypothesis is not decoration]
In Lean's mathematical library division is a total function: $0^{-1}$ is \emph{defined} to be
$0$. Without \lean{hR} the claimed inequality at $R=0$ would read $2\le0+0$, which is false, and
for $R<0$ it fails as well ($R=-1$ gives $2\le-2$). Lean would refuse the proof. When reading
any formal statement, first find the hypotheses and ask what each one excludes; a theorem
with a forgotten hypothesis is unprovable, and a theorem with an \emph{extra}, unsatisfiable
hypothesis is provable and empty.
\end{pitfallbox}

\subsection{What the inequality means, and with what confidence}

\Cref{thm:intro:amgm} is elementary mathematics, and its Lean form is Tier~A. Its
\emph{interpretation} is a different kind of statement. That $\sqrt{\ap}$ acts as a minimal
radius for a string on a circle is a result of the literature \TL{}
\source{Tong §8.3, ll.~11659--11663}. That the same mechanism matters in the early universe
is the proposal of string gas cosmology, which starts from the observation that strings at
small scale factor behave like strings at large scale factor and therefore act on a
cosmological background differently from point particles
\source{BW §III.B, ll.~720--737}; the same review describes scenarios in which compact
dimensions are driven towards the self-dual radius, and their fine-tuning problems \TL{}
\source{BW §I, ll.~160--176; §IV, ll.~1108--1137}. That $R_{\rm eff}$ itself is \emph{the}
physically measured length, or that the bound removes the initial singularity of cosmology,
is a proposal of this programme \TC{} and is examined as such in
\cref{ch:dualscale,ch:cosmology}. The general version for a $d$-dimensional torus with metric
$G$, $\Tr G+\Tr G^{-1}\ge2d$, is proved in \cref{ch:dualscale}; in Lean it is
\lean{dualScale\_ge}, in the same file as the circle bound \TA.

%=====================================================================================
\section{The programme of this book}\label{sec:intro:programme}

The circle is the first case of a pattern. Replace the circle by a $d$-dimensional torus: the
pair $(n,w)$ becomes a vector of $2d$ integers, the exchange $n\leftrightarrow w$ becomes the
group $\Odd{d}$ of integer matrices preserving the pairing $2\,n\!\cdot\! w$, and the radius
becomes a metric $G$ and a two-form $B$ assembled in the generalized metric $\HH$
(Parts~II and~III). Replace the torus by a K3 surface: the lattice of charges becomes the
even unimodular lattice of signature $(4,20)$, whose automorphisms again act as dualities
(Part~IV). Compactify on $\KT$ and add fluxes: integrality conditions and the number
$\chi/24$ appear (Part~V). Count states on K3: the coefficients organize into dimensions of
representations of the sporadic group $\Mtf$ (Part~VI). In each case the physics rests on a
piece of exact, discrete mathematics.

This observation defines the scope of the formalization. Its goal, in the words of the
project's workflow document, is a kernel-checked Lean~4 layer for the \emph{mathematical
skeleton} of string theory on $\KT$: charge lattices, their duality groups, the doubled
algebra, flux and tadpole counting, and the moonshine numerics
\source{docs/STREAM2\_WORKFLOW.md §1}. The same document states what is not a goal and is
never to be claimed: formalizing string theory as physics. Moduli-space quotients, the
classification of even unimodular lattices, Fourier--Mukai transforms and similar deep
results remain Tier~L unless a specific work package proves them. A student should read every
chapter with this division in mind: \emph{the physics is explained and sourced; the skeleton
is proved; the bridge between them is argued, never machine-checked.}

\begin{remark}[Older and newer libraries]\label{rem:intro:shadows}
The repository contains libraries of different vintage. The newest, \lean{DualScaleStream2},
works with real and integer matrices of arbitrary size and proves statements that are
recognizably the mathematical ones (the circle bound above is an example). Several older
libraries model physical quantities by integers or rationals and prove arithmetic
consequences. This book calls such results \emph{arithmetic shadows}\index{arithmetic shadow}
of the physical statement, says so each time, and explains what a faithful formalization
would need. \Cref{ch:architecture} maps the libraries.
\end{remark}

%=====================================================================================
\section{Three tiers of confidence}\label{sec:intro:tiers}

Every claim that matters in this book carries one of three badges.\index{epistemic tiers}
\begin{description}[leftmargin=2.2em,labelwidth=1.6em,itemsep=2pt,font=\normalfont]
\item[\TA] \emph{Kernel-checked.} A named declaration in the Lean~4 libraries of this
project, compiled without \texttt{sorry}. It may depend only on the three standard axioms of
Lean, namely \lean{propext}, \lean{Classical.choice} and \lean{Quot.sound}.
\item[\TL] \emph{Literature.} Established in a published source, cited with the line range
that was read; not formalized here.
\item[\TC] \emph{Conjecture.} An open conjecture of the field, or a proposal of this
programme.
\end{description}
Two rules govern the badges. First, a computer-algebra verification is a ``symbolic check'',
useful but not Tier~A. Second, and more important: \emph{the identification of a Lean object
with a physical system is never Tier~A.} The theorem of \cref{sec:intro:amgm} is about a real
number; the sentence ``the radius of a compact dimension, measured by strings, is bounded
below'' is Tier~L. The second rule is not a matter of taste. It follows from a simple
principle about how confidence propagates, and that principle has itself been formalized.

\subsection{The tier order and the ledger, in Lean}

The module \lean{DualScaleM24Formalization/Epistemic/Tier.lean} defines a type of tiers. It is
finer than the book's: it has five values, in increasing order of strength
\[
  \texttt{X}\;<\;\texttt{C}\;<\;\texttt{L}\;<\;\texttt{B}\;<\;\texttt{A},
\]
where \texttt{X} (exploratory: floating-point numerics, sampling, machine-generated text) and
\texttt{B} (a deterministic exact-arithmetic check accompanied by a negative control that
fails) are used in the project's internal ledgers. This book does not print \texttt{X}
material and files its symbolic checks in the running text, so three badges suffice. The
order is defined by sending each tier to a natural number and comparing
(\lean{Tier.toNat}, \lean{Tier.le}); the module proves that $\le$ is reflexive, transitive and
antisymmetric (\lean{tier\_refl}, \lean{tier\_trans}, \lean{tier\_antisymm}), that \texttt{A} is
the top and \texttt{X} the bottom element (\lean{tier\_A\_is\_maximal},
\lean{tier\_X\_is\_minimal}), and the small lemma on which everything else turns:
\begin{leancode}
theorem eq_A_of_A_le {s : Tier} (h : Tier.A ≤ s) : s = Tier.A := by ...
\end{leancode}
if something is at least as strong as \texttt{A}, it \emph{is} \texttt{A}.

The module \lean{DualScaleM24Formalization/Epistemic/Ledger.lean} models the bookkeeping. A
\emph{claim} is a tier together with the list of claims it cites; a \emph{ledger}%
\index{ledger (epistemic)} assigns a claim to each identifier of some index type $\iota$; a
ledger is \emph{sound} if no claim is filed above anything it directly cites.
\begin{leancode}
structure Claim (ι : Type) where
  tier : Tier
  supports : List ι

abbrev Ledger (ι : Type) := ι → Claim ι

def Sound {ι : Type} (L : Ledger ι) : Prop :=
  ∀ a b, b ∈ (L a).supports → (L a).tier ≤ (L b).tier
\end{leancode}
Soundness is a \emph{local} condition: it looks one citation deep. What one wants is a
\emph{global} guarantee, along chains of citations of any length. The relation
\lean{Depends L a b} (``$a$ rests on $b$, directly or through a chain'') is defined
inductively, and the main theorem states that the local condition implies the global one.

\begin{leanbox}[In Lean: confidence is a weakest-link quantity \TA]
\begin{leancode}
theorem tier_le_of_depends {ι : Type} {L : Ledger ι} (hL : Sound L) :
    ∀ {a b : ι}, Depends L a b → (L a).tier ≤ (L b).tier := by ...

theorem no_kernel_claim_rests_on_weaker {ι : Type} {L : Ledger ι}
    (hL : Sound L) {a b : ι} (ha : (L a).tier = Tier.A)
    (hab : Depends L a b) : (L b).tier = Tier.A := by ...
\end{leancode}
\textbf{Reading.} In a sound ledger, a claim is never filed above anything in its transitive
support; in particular everything a Tier-\texttt{A} claim rests on, at any depth, is
Tier~\texttt{A}. The contrapositive is \lean{not\_A\_of\_weak\_support}: one weaker link anywhere
in the chain forbids the badge \texttt{A}.

\textbf{Proof idea.} Induction on the chain. A chain of length one is the soundness
hypothesis; a longer chain is a first citation followed by a shorter chain, and the two
inequalities compose by \lean{tier\_trans}. The corollary rewrites the tier of $a$ as
\texttt{A} and applies \lean{eq\_A\_of\_A\_le}.

\textbf{What it does not say.} The theorem is about an abstract function
$\iota\to\lean{Claim}\ \iota$. It does not inspect this book. Whether the badges printed in
these pages form a sound ledger is an editorial discipline, which the theorem motivates but
cannot enforce.
\end{leanbox}

\begin{example}[A sound and an unsound ledger]\label{ex:intro:ledger}
The Lean module contains two test ledgers, drawn in \cref{fig:intro:ledger}. In
\lean{witnessSound}, claim $0$ is Tier~\texttt{A} and cites claim $1$, also Tier~\texttt{A};
claim $2$ is a free-standing conjecture. The only citation to check is $0\to1$, and
$\texttt{A}\le\texttt{A}$ holds: \lean{witnessSound\_sound} proves
\lean{Sound witnessSound}. In \lean{witnessUnsound}, claim $0$ is filed at \texttt{A} while
citing a Tier-\texttt{C} claim; in numbers, soundness would need $4\le1$. The theorem
\lean{witnessUnsound\_not\_sound} proves \lean{¬ Sound witnessUnsound}. The second test is as
important as the first: a definition of ``sound'' that every ledger satisfied would make the
main theorem true and useless, and the negative control rules this out.

Now apply the principle to this chapter. The physical sentence ``a string cannot resolve a
radius below $\sqrt{\ap}$'' cites (i) the inequality, Tier~A; (ii) the mass formula
\eqref{eq:intro:mass} and the extension of T-duality to the full theory, Tier~L; (iii) the
identification of the real number $R$ in Lean with a radius in string units, which no kernel
can check. By \lean{tier\_le\_of\_depends} the sentence can be filed at Tier~L at best. This is
the rule stated above, now seen as a consequence of soundness.
\end{example}

\begin{figure}[t]
\centering
\begin{tikzpicture}[node distance=14mm,
  cl/.style={draw,rounded corners=2pt,minimum width=11mm,minimum height=6.5mm,font=\small}]
  \begin{scope}
    \node[cl,draw=tierA,text=tierA] (a0) {$0$: A};
    \node[cl,draw=tierA,text=tierA,below=9mm of a0] (a1) {$1$: A};
    \node[cl,draw=tierC,text=tierC,right=8mm of a0] (a2) {$2$: C};
    \draw[-{Stealth}] (a0) -- node[right]{\scriptsize cites} (a1);
    \node[below=3mm of a1,font=\small] {sound};
  \end{scope}
  \begin{scope}[xshift=5.6cm]
    \node[cl,draw=tierA,text=tierA] (b0) {$0$: A};
    \node[cl,draw=tierC,text=tierC,below=9mm of b0] (b1) {$1$: C};
    \draw[-{Stealth}] (b0) -- node[right]{\scriptsize cites} (b1);
    \node[below=3mm of b1,font=\small] {unsound: $\texttt{A}\not\le\texttt{C}$};
  \end{scope}
  \begin{scope}[xshift=10.2cm]
    \node[cl,draw=tierL,text=tierL] (c0) {physical claim: L};
    \node[cl,draw=tierA,text=tierA,below left=9mm and -9mm of c0] (c1) {inequality: A};
    \node[cl,draw=tierL,text=tierL,below right=9mm and -9mm of c0] (c2) {spectrum: L};
    \draw[-{Stealth}] (c0) -- (c1);
    \draw[-{Stealth}] (c0) -- (c2);
    \node[below=13.5mm of c0,yshift=-6mm,font=\small] {sound: $\texttt{L}\le\texttt{A}$, $\texttt{L}\le\texttt{L}$};
  \end{scope}
\end{tikzpicture}
\caption{Ledgers as directed graphs; an arrow points from a claim to a claim it cites. Left
and middle: the two test ledgers of the Lean module (\cref{ex:intro:ledger}). Right: the way
this book files a physical statement that uses a kernel-checked inequality: at the tier of its
weakest support.}
\label{fig:intro:ledger}
\end{figure}

%=====================================================================================
\section{Why a proof assistant}\label{sec:intro:whylean}

A proof assistant\index{proof assistant} is a program that checks a proof down to the axioms
of a fixed logical foundation. In Lean~4\index{Lean 4} a small trusted component, the
\emph{kernel}\index{kernel (Lean)}, accepts a theorem only if every step reduces to the rules
of the underlying type theory; the tactics that produced the proof (\lean{ring},
\lean{linarith}, a language model, a student) need not be trusted, because their output is
re-checked. \Cref{ch:leanprimer} teaches enough Lean to read and write the statements in this
book. Three properties make the tool suited to this subject.

\emph{It removes a specific kind of error.} Much of the mathematics in later chapters is
finite but unforgiving: an $8\times8$ integer matrix must have determinant exactly $1$; the
squares of $26$ character degrees must add up to a group order; a block-matrix identity must
hold with every sign right. Such statements are rarely doubted and often mistyped. The
project's own record has an instructive case: a table of the dimensions of the irreducible
representations of $\Mtf$ entered in an early library had two entries duplicated and two
missing, so that the squares summed to $351\,061\,029$ instead of $|\Mtf|=244\,823\,040
=2^{10}\cdot3^3\cdot5\cdot7\cdot11\cdot23$. The table was replaced by the one in the literature
and is now guarded by a theorem, \lean{M24RepDim\_sum\_sq}, which states the sum rule
\source{docs/STREAM2\_WORKFLOW.md §6, ``Side findings''}; see \cref{ch:m24}.

\emph{It makes the hypotheses explicit}, as the pitfall box of \cref{sec:intro:amgm} showed.
In physics texts, conditions such as ``$G$ positive definite'', ``$B$ antisymmetric'' or
``$R>0$'' are often left to context; a formal statement cannot leave them out.

\emph{It separates the statement from the proof.} Once the kernel accepts a proof, the only
thing a reader has to examine is whether the \emph{statement} says what the text claims.
This is the one check that no kernel can perform, and the project's workflow reserves a
separate, recorded statement-review step for it \source{docs/STREAM2\_WORKFLOW.md §4, step S3}.
The same document fixes the acceptance rule used for every Tier-A badge in this book: the
library builds, no \texttt{sorry} (Lean's placeholder for a missing proof) remains, and the
command \lean{\#print axioms} reports only the three standard axioms; the producer of a proof
never grades it \source{docs/STREAM2\_WORKFLOW.md §2, items 1--2}.

\begin{pitfallbox}[Common pitfall: what a green build does not mean]
A kernel-checked theorem is true \emph{as stated, from its hypotheses}. It is not evidence
that the statement models nature, nor that a toy statement about integers captures a
statement about quantum fields, nor that the definitions are the intended ones. A library can
compile perfectly and prove only trivialities. The habit to acquire is to read a formal
statement the way one reads a contract: literally, and with attention to what is absent.
\end{pitfallbox}

%=====================================================================================
\section{A reader's guide}\label{sec:intro:guide}

The book has eight parts. Parts~I--VI build the physics and mathematics, Part~VII presents
the programme's own proposals with their tiers, and Part~VIII documents the formalization.
{\small
\begin{longtable}{@{}lp{0.29\textwidth}p{0.52\textwidth}@{}}
\toprule
Part & Title & What it does\\
\midrule
I & Foundations of the String (\cref{ch:intro}--\cref{ch:backgrounds}) & Classical string,
quantization and critical dimension, the conformal field theory needed later, background
fields and the low-energy action.\\
II & Compactification and T-Duality (\cref{ch:circle}--\cref{ch:dbranes}) & The circle in
full, Buscher rules, the Narain lattice, the group $\Odd{d}$, the two-torus, D-branes.\\
III & Double Field Theory (\cref{ch:doubled}--\cref{ch:dftaction}) & Doubled coordinates, the
strong constraint, the generalized metric $\HH$, Courant brackets, the action.\\
IV & K3 Surfaces and Lattices (\cref{ch:lattices}--\cref{ch:stringsk3}) & Even unimodular
lattices, $E_8$ and $U$; K3 surfaces, periods and Torelli; Kummer surfaces; the Mukai
lattice; strings on K3.\\
V & $\KT$ and Fluxes (\cref{ch:k3t2}--\cref{ch:swampland}) & Compactification on $\KT$,
tadpoles and $\chi/24$, flux superpotentials, swampland constraints.\\
VI & Mathieu Moonshine (\cref{ch:ellgenus}--\cref{ch:moonshine}) & The elliptic genus of K3,
the group $\Mtf$, and the observed decompositions.\\
VII & The Dual-Scale Theory (\cref{ch:dualscale}--\cref{ch:triad}) & The inequality of this
chapter in general form; string gas cosmology; the programme's conjectures, labelled as
such.\\
VIII & The Formalization (\cref{ch:leanprimer}--\cref{ch:open}) & A Lean primer, the
architecture of the libraries, an atlas of the theorems, the proof pipeline, open problems.\\
\bottomrule
\end{longtable}}

\paragraph{Routes through the book.} A physics student can read Parts~I--III in order, take
\cref{ch:lattices} as a mathematical interlude and continue. A mathematics student who
accepts the mass formula \eqref{eq:intro:mass} as a definition can start at \cref{ch:narain},
read Parts~III--IV and VI, and return to Part~I for motivation. A reader mainly interested in
formalization should read this chapter, then \cref{ch:leanprimer}, and then the section
``What the machine has checked'' of each chapter.

\paragraph{Conventions for every chapter.} Each chapter opens with the question it answers,
derives its results, works at least one example with numbers, and ends with the
machine-checked section, a summary, graded exercises (\(\star\) routine, \(\star\star\)
requires an idea, \(\star\star\star\) a small project; some are to be done in Lean) and
further reading. Sources are cited as \source{key §section, ll.~lines}; the line numbers
refer to the plain-text copies of the papers in the project repository, so that every
citation can be located mechanically.

%=====================================================================================
\section{What the machine has checked}\label{sec:intro:lean}

This chapter used two groups of declarations. The first is the circle bound of
\cref{sec:intro:amgm}, in \lean{DualScaleStream2.DualScale}. The file that contains it proves
the general torus bound first. The circle bound is then proved directly, as shown above, and
not by specialization; what links the two is the following identity.

\begin{leanbox}[In Lean: from the torus to the circle \TA]
\begin{leancode}
theorem dualScale_circle (R : ℝ) (hR : R ≠ 0) :
    dualScale (!![R ^ 2] : Matrix (Fin 1) (Fin 1) ℝ)
      = (R + R⁻¹) ^ 2 - 2 := by ...
\end{leancode}
\textbf{Reading.} Here \lean{dualScale G} is defined as the trace of the generalized metric at
$B=0$, and \lean{dualScale\_eq} shows it equals $\Tr G+\Tr G^{-1}$. For the $1\times1$ metric
$G=(R^2)$ --- a circle of radius $R$ in units $\ap=1$ --- this is $R^2+R^{-2}$, and the theorem
is the identity $R^2+R^{-2}=(R+R^{-1})^2-2$.

\textbf{Proof idea.} Compute the inverse of a $1\times1$ matrix explicitly, reduce both traces
to their single entry, clear denominators (\lean{field\_simp}) and finish with \lean{ring}.

\textbf{Why it matters.} With $d=1$ the torus bound \lean{dualScale\_ge} gives
$(R+R^{-1})^2-2\ge2$, that is $R+R^{-1}\ge2$ for $R>0$: the circle bound is the $d=1$ case of
a statement about all positive-definite metrics, taken up in \cref{ch:genmetric,ch:dualscale}.
\end{leanbox}

The second group is the tier order and the ledger of \cref{sec:intro:tiers}, in the namespace
\lean{SocrateAI.Epistemic}. The table collects what was cited, with the tier of the Lean
statement and the tier of the use made of it in the text.

{\small
\begin{longtable}{@{}p{0.41\textwidth}p{0.42\textwidth}c@{}}
\toprule
Declaration & Literal content & Tier\\
\midrule
\lean{circle\_effective\_scale\_ge\_two} & $2\le R+R^{-1}$ for real $R>0$ & \TA\\
\lean{dualScale\_circle} & $\Tr G+\Tr G^{-1}=(R+R^{-1})^2-2$ for $G=(R^2)$, $R\ne0$ & \TA\\
\lean{dualScale\_ge} (preview) & $2d\le\Tr G+\Tr G^{-1}$ for real positive-definite $G$ & \TA\\
\lean{tier\_refl}, \lean{tier\_trans}, \lean{tier\_antisymm} & the five tiers are partially
ordered (in fact linearly, through \lean{Tier.toNat}) & \TA\\
\lean{tier\_A\_is\_maximal}, \lean{tier\_X\_is\_minimal}, \lean{eq\_A\_of\_A\_le} & \texttt{A} is
the top element, \texttt{X} the bottom & \TA\\
\lean{tier\_le\_of\_depends} & sound ledger $\Rightarrow$ tier is monotone along citation
chains & \TA\\
\lean{no\_kernel\_claim\_rests\_on\_weaker}, \lean{not\_A\_of\_weak\_support} & support of an
\texttt{A} claim is all \texttt{A}; contrapositive & \TA\\
\lean{witnessSound\_sound}, \lean{witnessUnsound\_not\_sound} & positive and negative control
ledgers & \TA\\
\lean{M24RepDim\_sum\_sq} (mentioned) & $\sum_{i}(\dim\rho_i)^2=244\,823\,040$ for the $26$
tabulated dimensions & \TA\\
\midrule
T-duality of the circle spectrum \eqref{eq:intro:tduality} & Tong §8.3; formalized versions in
\cref{ch:tduality} & \TL\\
$\sqrt{\ap}$ as a minimal radius & Tong §8.3 & \TL\\
$R_{\rm eff}$ as the physical length; singularity avoidance & proposal of this programme &
\TC\\
\bottomrule
\end{longtable}}
For each theorem in the table the command \lean{\#print axioms} reports only axioms among
\lean{propext}, \lean{Classical.choice} and \lean{Quot.sound}.

\begin{summarybox}
\begin{itemize}[leftmargin=*,itemsep=1pt]
\item Perturbative gravity is an expansion in $E^2/M_{\rm pl}^2$ and is non-renormalizable; it
is an effective theory that lacks its short-distance degrees of freedom. String theory
proposes that these are strings, characterized by one scale, $\ell_s=\sqrt{\ap}$, with tension
$T=1/(2\pi\ap)$.
\item On a circle $X\sim X+2\pi R$ a particle has momentum $n/R$ only; a closed string also
has winding $w$ with energy $|w|R/\ap$. The spectrum \eqref{eq:intro:mass} is invariant under
$R\mapsto\ap/R$ with $n\leftrightarrow w$: T-duality. A string has two rulers, $R$ and $\ap/R$.
\item The T-duality-invariant effective scale satisfies $R+\ap/R\ge2\sqrt{\ap}$, with equality
only at the self-dual radius; the proof is $(R-\sqrt{\ap})^2/R\ge0$. In Lean:
\lean{circle\_effective\_scale\_ge\_two} (the case $\ap=1$).
\item Claims carry tiers: \TA{} kernel-checked, \TL{} literature, \TC{} conjecture. Confidence
is a weakest-link quantity (\lean{tier\_le\_of\_depends}); identifying a Lean object with a
physical system is never Tier~A.
\item A \texttt{leanbox} gives a verbatim statement, a literal reading, a proof idea and a
tier. Read the hypotheses first.
\end{itemize}
\end{summarybox}

%=====================================================================================
\section*{Exercises}
\addcontentsline{toc}{section}{Exercises}

\begin{exercise}[$\star$]\label{exr:intro:units}
Using $\hbar c\approx1.97\times10^{-14}$ GeV\,cm, convert to centimetres (a) the reduced
Planck mass \eqref{eq:intro:planck}; (b) a string scale one order of magnitude below
$10^{19}$ GeV. Express the tension $T=1/(2\pi\ap)$ of case (b) in GeV$^2$.
\end{exercise}

\begin{exercise}[$\star$]\label{exr:intro:table}
Extend the table of \cref{ex:intro:spectrum} by all states with $|n|,|w|\le2$ and
$N+\tilde N\le3$ allowed by level matching, at $R=3$ and at $R=\tfrac13$ (with $\ap=1$).
Verify that the two lists of masses coincide and exhibit the bijection between states.
\end{exercise}

\begin{exercise}[$\star$]\label{exr:intro:plpr}
With $p_{L}=n/R+wR/\ap$ and $p_{R}=n/R-wR/\ap$, verify that
$\tfrac12(p_L^2+p_R^2)=n^2/R^2+w^2R^2/\ap^2$ and $p_L^2-p_R^2=4nw/\ap$. Show that under
$R\mapsto\ap/R$, $n\leftrightarrow w$ one has $p_L\mapsto p_L$ and $p_R\mapsto-p_R$, and deduce
\eqref{eq:intro:tduality} without expanding the mass formula.
\end{exercise}

\begin{exercise}[$\star\star$]\label{exr:intro:massless}
Find all states of \eqref{eq:intro:mass} with $M^2=0$ (a) at a generic radius, (b) at
$R=\sqrt{\ap}$. You should find four additional solutions with $n,w\in\{\pm1\}$ in case (b).
Why can they exist only at the self-dual radius? (Their meaning, an enhanced gauge symmetry,
is the subject of \cref{ch:tduality}.)
\end{exercise}

\begin{exercise}[$\star\star$]\label{exr:intro:eq}
State and prove in Lean the equality case missing from the circle bound:
\begin{leancode}
example (R : ℝ) (hR : 0 < R) (h : R + R⁻¹ = 2) : R = 1 := by
  sorry
\end{leancode}
Hint: from the identity $R+R^{-1}-2=(R-1)^2/R$ obtain $(R-1)^2=0$; \lean{field\_simp} at
\lean{h} followed by \lean{nlinarith [sq\_nonneg (R - 1)]} is one route.
\end{exercise}

\begin{exercise}[$\star\star$]\label{exr:intro:alpha}
State in Lean the inequality with a general $\ap$, in a form that avoids square roots:
for real $R>0$ and $a>0$, $4a\le(R+a/R)^2$. Prove it on paper from \cref{thm:intro:amgm}, then
in Lean. Explain why the square-root-free form is equivalent to $R+a/R\ge2\sqrt a$ when
$R,a>0$. Lean's linter will report that the hypothesis $a>0$ is never used: why is the
square-root-free statement true for every real $a$, and why is the original one not?
\end{exercise}

\begin{exercise}[$\star\star$]\label{exr:intro:ledger}
In the ledger of \cref{fig:intro:ledger} (right), suppose an author refiles the physical claim
at Tier~\texttt{A}. Which instance of the condition \lean{Sound} fails? Write the four-element
ledger (the claim, the inequality, the spectrum, and a Tier-\texttt{C} interpretation cited by
nobody) as a Lean definition \lean{Ledger (Fin 4)} in the style of \lean{witnessSound}, and
prove it sound with the same three-line proof.
\end{exercise}

\begin{exercise}[$\star\star\star$]\label{exr:intro:invariants}
(a) Show that every function of $R>0$ invariant under $R\mapsto\ap/R$ can be written as a
function of $R_{\rm eff}=R+\ap/R$ alone. Hint: for $s\ge2\sqrt{\ap}$ the equation $R+\ap/R=s$
has exactly the two solutions $R$ and $\ap/R$. (b) Deduce that $\max(R,\ap/R)$, the size read
off from the lightest tower, is an increasing function of $R_{\rm eff}$, and find it explicitly.
(c) Discuss: does (a) single out $R_{\rm eff}$ as ``the'' physical length? Which tier does your
answer carry?
\end{exercise}

%=====================================================================================
\section*{Further reading}
\addcontentsline{toc}{section}{Further reading}

\begin{itemize}[leftmargin=*,itemsep=2pt]
\item Motivation for string theory and the problem of quantum gravity: \source{Tong §0,
ll.~369--680}. The tension, $\ap$ and the string length: \source{Tong §1.2, ll.~1090--1150}.
The string length as a minimum length in scattering: \source{Tong §6, ll.~8055--8072}.
\item The circle, winding and T-duality at the level used here:
\source{Tong §8.2--8.3, ll.~11345--11660}; continued in \cref{ch:circle,ch:tduality}.
\item The scope of target-space duality, from circles to curved backgrounds, and the history
of the word ``duality'': \source{GPR §1, ll.~207--330}.
\item Strings as a gas in the early universe and the role of T-duality there:
\source{BW §I, ll.~128--190; §III.B, ll.~720--737}; continued in \cref{ch:cosmology}.
\item The project's rules for what counts as checked, and the measured record of the proof
pipeline: \source{docs/STREAM2\_WORKFLOW.md §§1--2, 4, 6}; continued in
\cref{ch:pipeline}.
\end{itemize}
